\chapter{Quantum Singleton Bound}
\label{chap:quantum-singleton}
%%%%%%%%%%%%%%%%%%%%%%%%%%%%%%%%%%%%%%%%%%%%%%%%%%%%%%%%%%%%%%%%%%%%%%%%%%%%%%%
\section{Overview}

This chapter proves the \textbf{quantum Singleton bound}:
for a stabilizer code with logical dimension $k$ and distance $d$
on $n$ qudits over a prime field $\mathbb{F}_p$,
\[
k + 2(d-1) \;\le\; n.
\]
The proof uses the symplectic vector space $V = \mathbb{F}_p^n \times \mathbb{F}_p^n$
with the standard symplectic form, and the key idea of erasure correctability.

%%%%%%%%%%%%%%%%%%%%%%%%%%%%%%%%%%%%%%%%%%%%%%%%%%%%%%%%%%%%%%%%%%%%%%%%%%%%%%%
\section{Symplectic vector space}

\begin{definition}[Symplectic form]
\lean{sym_form}
\leanok
For $u = (x,z), v = (x',z') \in V = \mathbb{F}_p^n \times \mathbb{F}_p^n$,
\[
\omega(u,v) \;=\; \sum_{i=0}^{n-1}(x_i z'_i - z_i x'_i).
\]
\end{definition}

\begin{lemma}[Bilinearity]
\lean{sym_form_add_left}
\lean{sym_form_add_right}
\lean{sym_form_smul_left}
\lean{sym_form_smul_right}
\leanok
\uses{sym_form}
$\omega$ is bilinear: additive and scalar-homogeneous in each argument.
\end{lemma}

\begin{lemma}[Antisymmetry]
\lean{sym_form_swap}
\leanok
\uses{sym_form}
$\omega(u,v) = -\omega(v,u)$.
\end{lemma}

\begin{lemma}[Nondegeneracy]
\lean{sym_form_nondegenerate}
\leanok
\uses{sym_form}
If $\omega(u,v) = 0$ for all $v$, then $u = 0$.
\end{lemma}

\begin{definition}[Bundled bilinear form]
\lean{symB}
\leanok
\uses{sym_form}
$\omega$ packaged as a \lean{LinearMap.BilinForm}.
\end{definition}

%%%%%%%%%%%%%%%%%%%%%%%%%%%%%%%%%%%%%%%%%%%%%%%%%%%%%%%%%%%%%%%%%%%%%%%%%%%%%%%
\section{Support, weight, and isotropic subspaces}

\begin{definition}[Support and weight]
\lean{supp}
\lean{wt}
\leanok
The \emph{support} $\mathrm{supp}(v) \subseteq \mathrm{Fin}\,n$ consists of
coordinates where either component of $v$ is nonzero;
the \emph{weight} is $\mathrm{wt}(v) = |\mathrm{supp}(v)|$.
\end{definition}

\begin{lemma}[Weight bound]
\lean{wt_le_n}
\leanok
\uses{wt}
$\mathrm{wt}(v) \le n$ for all $v \in V$.
\end{lemma}

\begin{definition}[Support submodule]
\lean{V_sub}
\leanok
For $C \subseteq \mathrm{Fin}\,n$, the \emph{support submodule}
$V_C = \{v \in V \mid \mathrm{supp}(v) \subseteq C\}$.
\end{definition}

\begin{lemma}[Dimension of support submodule]
\lean{dim_V_sub}
\leanok
\uses{V_sub}
$\dim_{\mathbb{F}_p}(V_C) = 2|C|$.
\end{lemma}

\begin{definition}[Symplectic orthogonal complement]
\lean{sym_orth}
\leanok
\uses{sym_form}
For a submodule $S \le V$,
$S^{\perp_\omega} = \{v \in V \mid \forall s \in S,\; \omega(v,s) = 0\}$.
\end{definition}

\begin{lemma}[Dimension of the symplectic orthogonal]
\lean{finrank_sym_orth}
\leanok
\uses{sym_orth, sym_form_nondegenerate}
$\dim(S^{\perp_\omega}) = 2n - \dim(S)$.
\end{lemma}

\begin{definition}[Isotropic submodule]
\lean{IsIsotropic}
\leanok
\uses{sym_form}
$S$ is \emph{isotropic} if $S \le S^{\perp_\omega}$, i.e.\ $\omega(u,v) = 0$
for all $u,v \in S$.
\end{definition}

%%%%%%%%%%%%%%%%%%%%%%%%%%%%%%%%%%%%%%%%%%%%%%%%%%%%%%%%%%%%%%%%%%%%%%%%%%%%%%%
\section{Code parameters and erasure correctability}

\begin{definition}[Quantum code distance]
\lean{code_dist}
\leanok
\uses{wt, sym_orth, IsIsotropic}
$d(S) = \min\{\mathrm{wt}(v) \mid v \in S^{\perp_\omega} \setminus S,\; \mathrm{wt}(v) \neq 0\}$
(or $0$ if no such $v$ exists).
\end{definition}

\begin{lemma}[Distance is at most $n$]
\lean{code_dist_le_n}
\leanok
\uses{code_dist, wt_le_n}
$d(S) \le n$.
\end{lemma}

\begin{definition}[Erasure correctability]
\lean{correctable}
\leanok
\uses{sym_orth}
An erasure set $E \subseteq \mathrm{Fin}\,n$ is \emph{correctable} for $S$ if
every $v \in V_E \cap S^{\perp_\omega}$ is already in $S$.
\end{definition}

\begin{lemma}[Distance implies correctability]
\lean{dist_implies_correctable}
\leanok
\uses{correctable, code_dist}
If $|E| < d(S)$, then $E$ is correctable.
\end{lemma}

\begin{definition}[Logical dimension]
\lean{code_k}
\leanok
\uses{IsIsotropic}
$k(S) = n - \dim(S)$.
\end{definition}

%%%%%%%%%%%%%%%%%%%%%%%%%%%%%%%%%%%%%%%%%%%%%%%%%%%%%%%%%%%%%%%%%%%%%%%%%%%%%%%
\section{Key lemma: two disjoint correctable sets}

\begin{lemma}[Two disjoint correctable sets bound $k$]
\lean{two_disjoint_correctable_sets_bound_logical_dimension}
\leanok
\uses{correctable, code_k, dim_V_sub, finrank_sym_orth}
Let $S$ be isotropic.  If $A,B \subseteq \mathrm{Fin}\,n$ are disjoint and both correctable, then
\[
k(S) \;\le\; n - |A| - |B|.
\]
\end{lemma}

%%%%%%%%%%%%%%%%%%%%%%%%%%%%%%%%%%%%%%%%%%%%%%%%%%%%%%%%%%%%%%%%%%%%%%%%%%%%%%%
\section{Quantum Singleton bound}

\begin{theorem}[Quantum Singleton bound]
\lean{quantum_singleton_bound}
\leanok
\uses{two_disjoint_correctable_sets_bound_logical_dimension, dist_implies_correctable, code_dist_le_n}
For any isotropic submodule $S \le V$,
\[
k(S) + 2(d(S) - 1) \;\le\; n.
\]
\end{theorem}
